% =====================================================================
%  Part 2: Methodology
% =====================================================================
\section{Methodology}
\label{sec:method}

\subsection{Pipeline}

\begin{figure}[H]
  \centering
  \begin{tikzpicture}[node distance=4mm and 6mm, every node/.style={font=\small},
      box/.style={draw, rounded corners, align=center, minimum height=9mm, inner sep=3pt},
      >={Stealth}]
    \node[box] (raw) {Kaggle mirror\\images + titles};
    \node[box, right=of raw] (prep) {\texttt{prepare\_data}\\validity, clusters,\\split, resize};
    \node[box, right=of prep] (split) {committed\\\texttt{data/splits}};
    \node[box, right=of split] (train) {\texttt{train}\\UA, UB, F1, F2,\\ablations $\times$ 3 seeds};
    \node[box, right=of train] (eval) {\texttt{evaluate}\\test, subsets,\\cost};
    \node[box, below=of eval] (cmp) {\texttt{compare}\\decision rule,\\support/conflict};
    \node[box, left=of cmp] (ref) {K0 rule, T0,\\F0 late fusion};
    \node[box, left=of ref] (ana) {\texttt{analysis\_m3}, review:\\removal curves,\\Grad-CAM, taxonomy};
    \draw[->] (raw) -- (prep); \draw[->] (prep) -- (split); \draw[->] (split) -- (train);
    \draw[->] (train) -- (eval); \draw[->] (eval) -- (cmp); \draw[->] (eval) -- (ref); \draw[->] (ref) -- (cmp); \draw[->] (cmp) -- (ana);
  \end{tikzpicture}
  \caption{Pipeline. Every step is a script in the repository; every run stores its configuration, environment,
  log, history, selected checkpoint and test outputs.}
  \label{fig:pipeline}
\end{figure}

\subsection{Modality-Specific Encoders}
\label{sec:encoders}

\paragraph{Image encoder.}
We use ResNet-50~\cite{he2016resnet} with the torchvision ImageNet-1k weights \texttt{IMAGENET1K\_V2}, the same
backbone as P1 in Assignment~2, without its classification layer. A $224\times224$ crop gives the pooled feature
$\mathbf{u} = g_I(x^{I}) \in \R^{2048}$.

\paragraph{Text encoder.}
We use DistilBERT-base-uncased~\cite{sanh2019distilbert}, a 6-layer distilled BERT~\cite{devlin2019bert} with
hidden size $768$. Titles are HTML-unescaped and tokenised with its uncased WordPiece tokenizer (at most 64
tokens). Each batch is trimmed to the length of its longest title. The final hidden state of the
\texttt{[CLS]} token gives $\mathbf{v} = g_T(x^{T}) \in \R^{768}$.

\paragraph{Alignment.}
The two modalities are aligned at the item level through the shared page id. There is no word-to-region
annotation. F2 learns a soft token-level alignment with cross-modal attention between image regions and title
tokens (\cref{sec:f2}); the other models only align the two modalities at the item level.

\subsection{Models}
\label{sec:models}

\begin{table}[H]
  \centering\small
  \caption{Models. The parameter counts are printed by the code.}
  \label{tab:models}
  \begin{tabularx}{\linewidth}{@{}llYr@{}}
    \toprule
    ID & Modalities & Definition & Parameters \\
    \midrule
    K0 & text  & most-mentioned class in the title (lexicon, ties broken by the first mention); otherwise the most frequent train class & 0 \\
    T0 & text  & TF-IDF of word 1--2-grams and character 3--5-grams, multinomial logistic regression~\cite{pedregosa2011sklearn} & $13.6$\,M \\
    UA & image & $\operatorname{Linear}_{2048\to101}(\mathbf{u})$ & $23.71$\,M \\
    UB & text  & $\operatorname{Linear}_{768\to101}(\operatorname{Dropout}_{0.1}(\mathbf{v}))$ & $66.44$\,M \\
    F0 & both  & $\mathbf{p} = w\,\mathbf{p}_{\mathrm{UA}} + (1-w)\,\mathbf{p}_{\mathrm{UB}}$, with $w$ chosen on validation & $0$ extra \\
    F1 & both  & concat fusion, \cref{eq:f1} & $91.89$\,M \\
    F2 & both  & cross-modal transformer fusion with auxiliary losses and modality dropout, \cref{sec:f2} & $97.98$\,M \\
    \bottomrule
  \end{tabularx}
\end{table}

\paragraph{Late fusion (F0).}
F0 averages the softmax outputs of the UA and UB runs of the same seed. The weight is
$w^\star = \argmax_{w \in \{0, 0.05, \dots, 1\}} \mathrm{MacroF1}_{\mathrm{val}}(w)$. The selected weights are
$0.45$, $0.45$ and $0.50$ for seeds 42, 123 and 2026. F0 needs no training, so it is a reference that tells us how
much of the gain of a trained fusion comes from simply combining two good unimodal predictions.

\paragraph{Concat fusion (F1, intermediate fusion).}
Each modality is projected to $d=512$ and normalised. The two vectors are multiplied by availability masks
$m_I, m_T \in \{0,1\}$, concatenated and passed through an MLP:
\begin{equation}
  \label{eq:f1}
  \mathbf{h}_I = m_I \LN(W_I\mathbf{u}), \qquad
  \mathbf{h}_T = m_T \LN(W_T\mathbf{v}), \qquad
  f_{\mathrm{F1}} = W_2\,\operatorname{Dropout}_{0.1}\!\big(\operatorname{GELU}(W_1[\mathbf{h}_I;\mathbf{h}_T])\big),
\end{equation}
with $W_I \in \R^{512\times2048}$, $W_T \in \R^{512\times768}$, $W_1 \in \R^{512\times1024}$ and
$W_2 \in \R^{101\times512}$. The whole model, including both encoders, is trained end to end with one
cross-entropy loss on the fused logits.

\paragraph{Missing-modality handling.}
F1 zero-fills the projected vector of a missing modality ($m=0$). F0 falls back to the available unimodal model.
F1 is trained on complete pairs only, so zero-filling is outside its training distribution. F2 is trained
with missing modalities (modality dropout) and uses learned \texttt{[MISSING]} tokens instead.

\paragraph{Losses.}
UA, UB and F1 each use a single cross-entropy loss. F2 adds per-modality auxiliary losses (\cref{eq:f2loss}),
which are tested in ablation E1.

\subsection{Improved Fusion (F2)}
\label{sec:f2}

The M2 analysis showed that F1 under-uses the image. The title alone fits the training set, so the image branch
inside F1 ends up weaker than the standalone image model. F2 targets this failure with three changes, all fixed in
the proposal before any F2 run: token-level cross-modal attention, auxiliary unimodal losses, and modality dropout.

\begin{figure}[H]
  \centering
  \begin{tikzpicture}[every node/.style={font=\small},
      box/.style={draw, rounded corners, align=center, minimum height=8mm, inner sep=3pt},
      aux/.style={draw, rounded corners, align=center, inner sep=2pt, font=\scriptsize, dashed},
      >={Stealth}]
    \node[box] (img) at (0,0) {image\\$224\times224$};
    \node[box] (rn)  at (3.0,0) {ResNet-50 layer4\\$7\times7\times2048$};
    \node[box] (pi)  at (6.4,0) {$W_I\,\cdot + \mathbf{p}_i + \mathbf{t}_I$\\49 tokens $\times 512$};
    \node[box] (txt) at (0,-2.2) {title\\$\le 64$ tokens};
    \node[box] (db)  at (3.0,-2.2) {DistilBERT\\$T\times768$};
    \node[box] (pt)  at (6.4,-2.2) {$W_T\,\cdot + \mathbf{t}_T$\\$T$ tokens $\times 512$};
    \node[box, fill=gray!8, minimum height=20mm] (tf) at (9.6,-1.1) {2 pre-LN\\transformer\\layers, 8 heads};
    \node[box] (fus) at (9.6,0.65) {\texttt{[FUS]}};
    \node[box] (head) at (12.3,-1.1) {LN $\to$ 101\\$\mathcal{L}_{\mathrm{fused}}$};
    \node[aux] (ai) at (3.0,1.25) {aux head on pooled $\mathbf{u}$: $\lambda\,\mathcal{L}_I$};
    \node[aux] (at) at (3.0,-3.45) {aux head on \texttt{[CLS]} $\mathbf{v}$: $\lambda\,\mathcal{L}_T$};
    \draw[->] (img) -- (rn); \draw[->] (txt) -- (db);
    \draw[->] (rn) -- (pi); \draw[->] (db) -- (pt);
    \draw[->] (pi.east) -- (tf.west |- pi.east); \draw[->] (pt.east) -- (tf.west |- pt.east);
    \draw[->] (fus) -- (tf); \draw[->] (tf) -- (head);
    \draw[->, dashed] (rn) -- (ai); \draw[->, dashed] (db) -- (at);
  \end{tikzpicture}
  \caption{F2. Image regions and title tokens attend to each other and to a learned \texttt{[FUS]} token, whose final
  state is classified. Dashed: auxiliary unimodal heads. During training, one modality of $15\%$ of the items is
  replaced by a learned \texttt{[MISSING]} token.}
  \label{fig:f2}
\end{figure}

\paragraph{Tokens.}
The image gives the ResNet-50 \texttt{layer4} map, i.e.\ $49$ region features $\mathbf{f}_i\in\R^{2048}$. The title
gives the DistilBERT hidden states $\mathbf{h}_j\in\R^{768}$, which already contain word positions. Both are projected
to $d=512$:
\begin{equation}
  \mathbf{z}^I_i = W_I\mathbf{f}_i + \mathbf{p}_i + \mathbf{t}_I, \qquad
  \mathbf{z}^T_j = W_T\mathbf{h}_j + \mathbf{t}_T, \qquad
  Z^{(0)} = [\,\mathbf{z}_{\mathrm{FUS}};\ \mathbf{z}^I_1,\dots,\mathbf{z}^I_{49};\ \mathbf{z}^T_1,\dots,\mathbf{z}^T_{T}\,],
\end{equation}
where $\mathbf{p}_i$ is a learned position embedding of the $7\times7$ grid, $\mathbf{t}_I,\mathbf{t}_T$ are learned
type embeddings and $\mathbf{z}_{\mathrm{FUS}}$ is a learned fusion token.

\paragraph{Fusion.}
Two pre-LN transformer layers~\cite{vaswani2017attention} ($8$ heads, FFN $2048$, dropout $0.1$) run self-attention
over the whole sequence. Padding tokens are masked. Every image region can therefore attend to every title word
and vice versa. The prediction is $f_{\mathrm{F2}} = W_o\LN(\mathbf{z}^{(2)}_{\mathrm{FUS}})$. A single fusion token
that gathers both modalities is close to the bottleneck tokens of Nagrani et al.~\cite{nagrani2021bottlenecks}.
Unlike F1, this design can weigh individual regions and words differently for every item. The \texttt{[FUS]}
attention also shows how much the model relies on each modality (\cref{sec:qualitative}).

\paragraph{Auxiliary unimodal losses.}
Two linear heads classify the pooled image feature $\mathbf{u}\in\R^{2048}$ and the text \texttt{[CLS]} state
$\mathbf{v}\in\R^{768}$:
\begin{equation}
  \label{eq:f2loss}
  \mathcal{L} = \CE\big(f_{\mathrm{F2}}, y\big) + \lambda\,\big[\CE(a_I(\mathbf{u}), y) + \CE(a_T(\mathbf{v}), y)\big],
  \qquad \lambda = 0.5 .
\end{equation}
Each encoder therefore has to stay discriminative on its own, so the image encoder keeps learning even when the
title already fits the training labels~\cite{wang2020makes}. The auxiliary heads are not used for the prediction.

\paragraph{Modality dropout and missing modalities.}
During training, each complete pair loses one modality with probability $p=0.15$ (image or title, $50/50$, never
both), in the spirit of ModDrop~\cite{neverova2016moddrop}. All tokens of the dropped modality are replaced by one
learned \texttt{[MISSING\_IMG]} or \texttt{[MISSING\_TXT]} token; the other positions of that modality are masked.
The fusion therefore learns to classify from one modality alone. At test time, a missing modality is handled in the
same way, with no special case. The auxiliary losses always see both encoders.

\paragraph{Training schedule.}
In epoch~1 both encoders are frozen and only the new layers learn ($5\cdot10^{-4}$). The randomly initialised fusion
layers then do not send noisy gradients into the pretrained encoders. From epoch~2, the whole model is fine-tuned
with the same learning rates as F1. The warm-up covers $15\%$ of the steps, so the encoders also start with a
warm-up after they are unfrozen. The model has $97.98$\,M parameters: $8.11$\,M new layers on top of the two
encoders. It needs $5.06$\,GMACs per item, against $4.73$ for F1.

\subsection{Ablations and Analyses}
\label{sec:ablations}

\begin{table}[H]
  \centering\small
  \caption{Controlled experiments. Each changes one factor of F2, or of the input, and keeps everything else, including
  the seeds and the split.}
  \label{tab:ablations}
  \begin{tabularx}{\linewidth}{@{}lYY@{}}
    \toprule
    ID & Change & Question \\
    \midrule
    E1 & $\lambda = 0$. The auxiliary heads are trained on detached features, so they still measure each encoder without changing it & Do the auxiliary losses keep the image branch alive? \\
    E2 & every class-name mention (the lexicon of Part~1) is replaced by \texttt{[MASK]} in the train, validation and test titles; UB, T0, F0, F1 and F2 are retrained & How much of the result depends on the label leaking into the title, and which fusion uses the image best when it does not? \\
    E4 & $p = 0$ (no modality dropout) & Does the modality dropout help with missing modalities, and does it cost anything on complete pairs? \\
    E-mod & at test time, the title (or the image) is removed from a random $0$, $25$, $50$, $75$ or $100\%$ of the complete pairs, without retraining & How fast does each model degrade when a modality goes missing? \\
    \bottomrule
  \end{tabularx}
\end{table}

\paragraph{Notes on the ablations.}
\begin{itemize}
  \item \textbf{E2 masking.} It changes $79.8\%$ of the training titles ($50{,}092$ of $62{,}771$). After masking, the
        lexicon finds no class name in any title. The mention subsets of \cref{sec:main} still use the original
        titles.
  \item \textbf{E-mod fallbacks.} For a removed modality, F1 zero-fills it and F2 uses its \texttt{[MISSING]} token.
        F0 falls back to the other unimodal model. UA and UB predict the most frequent training class.
\end{itemize}

\paragraph{Analyses.}
\begin{itemize}
  \item \textbf{Support/conflict groups} over three seeds. The groups come from the UA and UB runs of the same seed.
  \item \textbf{Temperature scaling}~\cite{guo2017calibration}: one temperature per run, fitted on validation.
  \item \textbf{Qualitative examples} for each group:
        \begin{itemize}
          \item a Grad-CAM map~\cite{selvaraju2017gradcam} of F2's image encoder for the predicted class;
          \item the \texttt{[FUS]} attention of the last fusion layer, averaged over heads.
        \end{itemize}
  \item \textbf{Error taxonomy} from a manual review of $100$ conflict items, 50 Conflict-I and 50 Conflict-T,
        drawn at random from seed 42. Each reviewer chose the main reason why the failing unimodal model is wrong:
        \begin{itemize}
          \item image unrelated or visually similar (for the image model);
          \item title generic or title naming another dish (for the text model);
          \item wrong label or other (for both).
        \end{itemize}
        Phan Bá Thanh reviewed items 1--50 and Trần Công Hoàng Phước items 51--100. Hard cases were discussed
        together.
\end{itemize}

\subsection{Pretraining and Fine-Tuning}
\label{sec:training}

Both encoders start from public pretrained weights and are fully fine-tuned with discriminative learning rates.
\Cref{tab:setup} lists the setup; only the noted values differ between models.

\begin{table}[H]
  \centering\small
  \caption{Training setup.}
  \label{tab:setup}
  \begin{tabularx}{\linewidth}{@{}lY@{}}
    \toprule
    Image transforms & train: \texttt{RandomResizedCrop(224, scale 0.35--1)}, horizontal flip, colour jitter $0.1$; evaluation: \texttt{Resize(256)} + \texttt{CenterCrop(224)} \\
    Optimiser & AdamW~\cite{loshchilov2019adamw}. Learning rates: image encoder $10^{-4}$, text encoder $3\cdot10^{-5}$, new layers $10^{-3}$ (UA, UB) or $5\cdot10^{-4}$ (F1, F2). Weight decay $0.05$ (image encoder and heads), $0.01$ (text encoder); none on biases and normalisation parameters \\
    Schedule & linear warm-up over $5\%$ of the steps, then cosine decay to $1\%$ \\
    Epochs, batch & UA and F1: 10 epochs; UB: 5 epochs; batch size 64; bf16 autocast; gradient clipping at $1.0$ \\
    F2 and its ablations & as F1, plus: epoch 1 with frozen encoders, warm-up over $15\%$ of the steps, head learning rate $5\cdot10^{-4}$ \\
    Checkpoint & the epoch with the best validation macro-F1 on complete pairs; early stopping after 3 epochs without improvement \\
    Seeds & split seed $42$ (fixed); training seeds $42$, $123$, $2026$ \\
    \bottomrule
  \end{tabularx}
\end{table}

\subsection{Evaluation Protocol}
\label{sec:protocol}

\paragraph{Metrics.}
\begin{itemize}
  \item Macro-F1 (primary), top-1 and top-5 accuracy.
  \item Expected calibration error with $15$ equal-width bins~\cite{guo2017calibration},
        $\mathrm{ECE}=\sum_b \frac{|B_b|}{N}\,|\mathrm{acc}(B_b)-\mathrm{conf}(B_b)|$.
  \item Per-class F1 and the confusion matrix.
  \item Compute cost: parameters, training time, throughput at batch 128 and latency at batch 1.
\end{itemize}

\paragraph{Decision rule.}
The rule is the same as in A1 and A2 and was registered in the proposal. Model~A is declared better than model~B
only if both conditions hold:
\begin{enumerate}
  \item the difference in mean test macro-F1 over the three seeds exceeds the larger of the two seed standard deviations;
  \item a McNemar test~\cite{mcnemar1947} on the seed-42 predictions rejects equal error rates at $p<0.05$.
\end{enumerate}

\paragraph{Test views.}
Every model is evaluated on:
\begin{itemize}
  \item the complete test pairs ($22{,}594$);
  \item the dedup subset;
  \item the title-unseen subset;
  \item titles that do or do not mention the true class;
  \item the natural missing-text set ($121$ items).
\end{itemize}
The fusion models are also tested with each modality removed.

\subsection{Hardware and Software}
\label{sec:hardware}

The Kaggle GPU quota ran out during M2, so all training runs (M2 and M3) used a lab workstation:
\begin{itemize}
  \item \textbf{GPU:} NVIDIA RTX A6000 (48\,GB), shared with other users' inference servers.
  \item \textbf{Machine:} 64 CPU cores, Ubuntu 24.04.
  \item \textbf{Software:} Python 3.13.13, PyTorch 2.12.0+cu130, torchvision 0.27.0, Transformers 5.23.0, CUDA 13.0.
\end{itemize}
The runner (\texttt{scripts/remote/}) runs inside \texttt{tmux}. It picks the GPU with the most free memory
(32\,GB free at the start), logs the GPU usage every five minutes, and resumes from the last checkpoint after a
crash; no crash occurred in either milestone. Data preparation ran on a Kaggle CPU notebook with 4 vCPU
(26 minutes). The M3 runs used the same machine, software and $16$ data-loading workers. T0 and the analyses ran
on the CPU of the same machine.
